\subsection{Extension model architecture}
\label{sec:e_architecture}
The extension architecture is based on the baseline architecture of \autoref{sec:architecture}, with two modifications. Below discusses the reasoning behind each modification and how it is implemented. 
\subsubsection{Splitting the cue into parts}
\label{sec:splitting_cue}
The baseline architecture passes the network the TF-Map cue of \cite{zhang2025multi} and includes the energy-recovery projection (Equation~\ref{eq:tfmap-energy}). By Equation~\ref{eq:tfmap-unit}, every mixture frame is its loudness times its shape,
  $|X_t| = \big\lVert |X_t| \big\rVert\,\tilde{x}_t$. Substituting this into $m_t$ of Equation~\ref{eq:tfmap-energy} shows what the baseline's
  single number contains:
  \begin{equation}
    m_t = \big\langle\, |X_t|,\ u_t \,\big\rangle
        = \big\lVert |X_t| \big\rVert \, \big\langle \tilde{x}_t,\ u_t \big\rangle
        = \big\lVert |X_t| \big\rVert \, c_t ,
    \label{eq:matched-level}
  \end{equation}
  where 
  \begin{equation}
    c_t = \cos\theta_t = \big\langle \tilde{x}_t,\ u_t \big\rangle
    \label{eq:cue-similarity}
  \end{equation}
is the cosine similarity of Equation~\ref{eq:tfmap-cos} which is measured with the blended template of $u_t$ instead of an enrollment frame. The problem with $m_t$ in the baseline architecture is that once multiplying the loudness and the similarity into a scalar, the network cannot tell distinguish between a frame that is very loud but not similar to the target and a frame that is very similar but not loud. The idea is to pass the parts instead of the product. Two parts are already defined: the template direction $u_t$ of Equation~\ref{eq:tfmap-energy} and the similarity $c_t$ of Equation~\ref{eq:cue-similarity}. The third part is new and is the residual
\begin{equation}
  r_t = |X_t| - m_t u_t
  \label{eq:cue-residual}
\end{equation}
which is the part of the mixture frame that the template of the target does not explain. A property of the residual is that it is orthogonal to the template, shown by:
\begin{equation}
  \big\langle r_t, u_t \big\rangle
  = \big\langle |X_t| - m_t u_t, u_t \big\rangle
  = \big\langle |X_t|, u_t \big\rangle - m_t \big\langle u_t, u_t \big\rangle
  = m_t - m_t = 0.
  \label{eq:cue-residual-orthogonal}
\end{equation}
where the last equality uses the fact that $u_t$ is a unit vector. The orthogonality thus allows a frame to be split exactly into the equation:
\begin{equation}
  |X_t| = m_t u_t + r_t
  \label{eq:cue-split}
\end{equation}
Therefore we can rebuild $\phi_t$ exactly as:
\begin{equation}
  \phi_t = c_t\big\lVert |X_t| \big\rVert u_t
  \label{eq:cue-rebuild}
\end{equation}
The choice to split was this work's own design choice, but was a suggestion made by AI model Claude \cite{anthropic2026claude} to diagnose the baseline's inability to learn the target's voice effectively.

\subsubsection{Speaker context embedding}
\label{sec:speaker_context_embedding}
When considering Section~\ref{sec:splitting_cue} the cue is still only described by the enrollments spectra. Target speakers that have similar pitch and vocal characteristics thus produce similar templates. Measured on the baseline architecture, the extractor was able to select the right voice \SI{71.8}{\percent} of the time when two speakers were of different gender, but only \SI{52.6}{\percent} when they were the same gender. The idea for this extension is to incorporate further target speaker identity. The second extension is to add a speaker embedding for each target speaker. Each enrollment clip is passed through a frozen speaker recognition network, ECAPA-TDNN \cite{desplanques2020ecapa} using the SpeechBrain checkpoint trained on VoxCeleb \cite{ravanelli2021speechbrain}, to produce a 192-dimensional embedding. The embedding is then passed through a small learned layer to produce a 128-dimensional vector to be used in the extractor. ECAPA itself is large at 20.77 million parameters and hence would cause storage concerns if it were included in the extractor. However, the ECAPA network can be run once per enrollment clip and the resulting embedding can be stored for later use. This allows for the extractor to be served locally on a device without needing to have the full ECAPA network, only the stored embedding for the target speaker.

The WeSep toolkit \cite{wang2024wesep}, an established open-source implementation of the baseline architecture, supports two uses of the embedding. The first (which is the one used in this work) is to consider utterance-level embedding, and the second is to consider contextual embedding of Zhang et al. \cite{zhang2025multi}. Although WeSep supports the contextual embedding, it was discarded as an option for this work because it recomputes the embedding for every frame of the mixture, which is computationally expensive. 
